\chapter*{Part Synthesis: Modalities}
\addcontentsline{toc}{chapter}{Part Synthesis: Modalities}

\begin{whychaptermatters}
This part made the book's general workflow modality-specific. Vision sharpened invariance, language sharpened context and grounding, and sequential modeling sharpened prediction-time legality. The common lesson is that different data types still demand explicit audits before architecture choice can be trusted.
\end{whychaptermatters}

\begin{chaptersummary}
\item \textbf{Question this part taught.} What task-specific structure must be preserved, what nuisance variation should be ignored, and what evidence makes a modality-specific claim honest under deployment conditions?
\item \textbf{Artifact chain this part added.} Vision design and stress-test card $\rightarrow$ NLP context-and-grounding card $\rightarrow$ sequential design card.
\item \textbf{What a student can now do.} Match invariance, context, grounding, and prediction-time audits to the modality instead of treating pixels, tokens, and sequences as interchangeable inputs.
\end{chaptersummary}

